%------------------------------------------------------------------------------%
\documentclass[fleqn,utf8,aspectratio=169,12pt]{beamer}

\usepackage{gaal}

\title{Propriedades dos Determinantes}

%------------------------------------------------------------------------------%
\begin{document}

\addtitle

%------------------------------------------------------------------------------%
\section{Propriedades dos Determinantes}

%------------------------------------------------------------------------------%
\begin{frame}{Propriedade 1 -- Matriz triangular}
Se $A$ é uma matriz triangular superior ou inferior
\[
  \pause
  A =
  \begin{pmatrix}[cccc]
    a_{11} & a_{12} & \cdots & a_{1n} \\
    0      & a_{22} & \cdots & a_{2n} \\
    \vdots & \vdots & \ddots & \vdots \\
    0      & 0      & \cdots & a_{nn}
  \end{pmatrix}
  \pause
  \qquad\text{ ou }\qquad
  A =
  \begin{pmatrix}[cccc]
    a_{11} & 0      & \cdots & 0      \\
    a_{21} & a_{22} & \cdots & 0      \\
    \vdots & \vdots & \ddots & \vdots \\
    a_{n1} & a_{n2} & \cdots & a_{nn}
  \end{pmatrix}
\]

\pause
então
\[
  \det(A) = a_{11}\; a_{22}\; \cdots\; a_{nn}
\]
\end{frame}

%------------------------------------------------------------------------------%
%------------------------------------------------------------------------------%
\begin{question}
  Calcule o determinante da matriz
  \[
    A =
    \begin{pmatrix}
      2 & 3 & 1  \\
      0 & 4 & 5  \\
      0 & 0 & -2
    \end{pmatrix}
  \]
\end{question}

%------------------------------------------------------------------------------%
\begin{resolution}{Solução}
  \onslide<+->{\[
    A =
    \begin{pmatrix}
      \emph{2} & 3        & 1         \\
      \z       & \emph{4} & 5         \\
      \z       & \z       & \emph{-2}
    \end{pmatrix}
  \]}

  \begin{align*}
    \onslide<+->{\det(A)}
    \onslide<+->{& = a_{11} \; a_{22} \; a_{22} \\}
    \onslide<+->{& = 2 \times 4 \times (-2) \\}
    \onslide<+->{& = -16}
  \end{align*}
\end{resolution}

%------------------------------------------------------------------------------%
\endinput


%------------------------------------------------------------------------------%
\begin{frame}{Propriedade 2 -- Matriz diagonal}
  Se $A$ é diagonal
  \[
    A =
    \begin{pmatrix}[cccc]
      a_{11} & 0      & \cdots & 0      \\
      0      & a_{22} & \cdots & 0      \\
      \vdots & \vdots & \ddots & \vdots \\
      0      & 0      & \cdots & a_{nn}
    \end{pmatrix}
  \]
  \pause
  então
  \[
    \det(A) = a_{11} \; a_{22} \; \cdots \; a_{nn}
  \]

  \pause
  Em particular
  \[
    \det(I_n) = 1
  \]
\end{frame}

%------------------------------------------------------------------------------%
\begin{frame}{Propriedade 3 -- Troca de linhas}
  Quando duas linhas de uma matriz são trocadas
  \[
    L_i \leftrightarrow L_j
  \]
  o determinante muda de sinal
  \[
  \det(B) = -\det(A)
  \]
\end{frame}

%------------------------------------------------------------------------------%
%------------------------------------------------------------------------------%
\begin{question}
  Calcule o determinante da matriz
  \[
    A =
    \begin{pmatrix}
      1 & 2 \\
      3 & 4
    \end{pmatrix}
  \]
  Crie a matriz $B$ trocando as linhas \(L_1 \leftrightarrow L_2\)

  Calcule \(\det(B)\)
\end{question}

%------------------------------------------------------------------------------%
\begin{resolution}{Solução}
\begin{columns}
\column{0.4\linewidth}
  \onslide<+->{\[
    A=
    \begin{pmatrix}
      1 & 2 \\
      3 & 4
    \end{pmatrix}
  \]}
  \begin{align*}
    \onslide<+->{\det(A)}
    \onslide<+->{& = 1 \times 4 - 2 \times 3 \\}
    \onslide<+->{& = 4 - 6 \\}
    \onslide<+->{& = -2}
  \end{align*}
\column{0.4\linewidth}
  \onslide<+->{\[
    B =
    \begin{pmatrix}
      3 & 4 \\
      1 & 2
    \end{pmatrix}
  \]}
  \begin{align*}
    \onslide<+->{\det(B)}
    \onslide<+->{& = 3 \times 2 - 4 \times 1 \\}
    \onslide<+->{& = 6 - 4 \\}
    \onslide<+->{& = 2 \\}
    \onslide<+->{& = -\det(A)}
  \end{align*}
\end{columns}
\end{resolution}

%------------------------------------------------------------------------------%
\endinput


%------------------------------------------------------------------------------%
\begin{frame}{Propriedade 4 -- Multiplicação de uma linha}
  Quando todos os elementos de uma linha são multiplicados por $k$
  \pause
  \[
    L_i \leftarrow k L_i
  \]

  \pause
  o determinante também é multiplicado por $k$
  \pause
  \[
    \det(B) = k\det(A)
  \]
\end{frame}

%------------------------------------------------------------------------------%
\begin{frame}{Propriedade 4 -- Multiplicação de uma linha}
  Note que se multiplicarmos todas as linhas por $k$

  \pause
  o determinante é multiplicado por $k^n$
\end{frame}

%------------------------------------------------------------------------------%
%------------------------------------------------------------------------------%
\begin{question}
  Calcule o determinante da matriz
  \[
    A =
    \begin{pmatrix}
      1 & 2 \\
      3 & 4
    \end{pmatrix}
  \]
  Crie a matriz $B$ multiplicando a primeira linha por 3

  Calcule \(\det(B)\)
\end{question}

%------------------------------------------------------------------------------%
\begin{resolution}{Solução}
\begin{columns}
\column{0.35\linewidth}
  \onslide<+->{\[
    A=
    \begin{pmatrix}
      1 & 2 \\
      3 & 4
    \end{pmatrix}
  \]}
  \begin{align*}
    \onslide<+->{\det(A)}
    \onslide<+->{& = 1 \times 4 - 2 \times 3 \\}
    \onslide<+->{& = 4 - 6 \\}
    \onslide<+->{& = -2}
  \end{align*}
\column{0.5\linewidth}
  \onslide<+->{\[
    B =
    \begin{pmatrix}
      3 & 6 \\
      3 & 4
    \end{pmatrix}
  \]}
  \begin{align*}
    \onslide<+->{\det(B)}
    \onslide<+->{& = 3 \times 4 - 6 \times 3 \\}
    \onslide<+->{& = 12 - 18 \\}
    \onslide<+->{& = -6 \\}
    \onslide<+->{& = 3\det(A)}
  \end{align*}
\end{columns}
\end{resolution}

%------------------------------------------------------------------------------%
\endinput


%------------------------------------------------------------------------------%
\begin{frame}{Propriedade 5 -- Adição de múltiplo de uma linha}
  Adicionar um múltiplo de uma linha
  \[
    L_i \leftarrow L_i + k L_j
    \qquad i\neq j
  \]

  \pause
  \emph{não altera o determinante}
  \[
    \det(B) = \det(A)
  \]
\end{frame}

%------------------------------------------------------------------------------%
%------------------------------------------------------------------------------%
\begin{question}
  Calcule o determinante da matriz por escalonamento
  \[
    A =
    \begin{pmatrix}
      1 & 2 & 3 \\
      2 & 5 & 7 \\
      1 & 1 & 2
    \end{pmatrix}
  \]
\end{question}

%------------------------------------------------------------------------------%
\begin{resolution}{Solução}
\end{resolution}

%------------------------------------------------------------------------------%
\endinput

Aplicamos
\[
L_2\leftarrow L_2-2L_1
\]
e
\[
L_3\leftarrow L_3-L_1.
\]

Essas operações não alteram o determinante
\[
\det(A)=
\begin{vmatrix}
1&2&3\\
0&1&1\\
0&-1&-1
\end{vmatrix}.
\]

Agora aplicamos
\[
L_3\leftarrow L_3+L_2.
\]

Obtemos
\[
\det(A)=
\begin{vmatrix}
1&2&3\\
0&1&1\\
0&0&0
\end{vmatrix}.
\]


Como a matriz é triangular,
\[
\det(A)=1\cdot1\cdot0=0.
\]



%------------------------------------------------------------------------------%
\begin{frame}{Propriedade 6 -- Linha nula}
  Se uma matriz possui uma linha ou coluna composta somente por zeros
  \pause
  \[
    \det(A) = 0
  \]

  \pause
  \bigskip
  Consequência do Teorema de Laplace
\end{frame}

%------------------------------------------------------------------------------%
\begin{frame}{Propriedade 7 -- Linhas iguais}
  Se duas linhas de uma matriz são iguais
  \pause
  \[
    \det(A) = 0
  \]
\end{frame}

%------------------------------------------------------------------------------%
\begin{frame}{Propriedade 7 -- Linhas iguais}
  Se trocarmos duas linhas o determinante muda de sinal
  \[
    \pause
    \det(B) = - \det(A)
  \]
  \pause
  Mas como as linhas são iguais, a matriz não muda \(B=A\)
  \[
    \pause
    \det(A) = - \det(A)
  \]
  \pause
  Logo
  \[
    \det(A) = 0
  \]
\end{frame}

%------------------------------------------------------------------------------%
\begin{frame}{Propriedade 8 -- Linhas proporcionais}
  Se duas linhas de uma matriz são proporcionais
  \[
    \det(A) = 0
  \]
\end{frame}

%------------------------------------------------------------------------------%
%------------------------------------------------------------------------------%
\begin{question}
  Calcule o determinante da matriz
  \[
    A =
    \begin{pmatrix}
      1 & 2 & 3 \\
      2 & 4 & 6 \\
      1 & 0 & 1
    \end{pmatrix}
  \]
\end{question}

%------------------------------------------------------------------------------%
\begin{resolution}{Solução}
  A segunda linha é duas vezes a primeira
  \[
    \det(A) = 0
  \]
\end{resolution}

%------------------------------------------------------------------------------%
\endinput

%------------------------------------------------------------------------------%
\begin{frame}{Propriedade 9 -- Transposta}
  O determinante de uma matriz é igual ao determinante de sua transposta
  \[
    \det\left(A^t\right) = \det(A)
  \]

  \pause
  As propriedades nas linhas também valem para as colunas
\end{frame}

%------------------------------------------------------------------------------%
\begin{frame}{Propriedade 10 -- Produto de matrizes}
  Se $A$ e $B$ são matrizes quadradas de mesma ordem
  \[
    \det(AB) = \det(A) \det(B)
  \]
\end{frame}

%------------------------------------------------------------------------------%
\begin{frame}{Propriedade 10 -- Produto de matrizes}
  Se $A$ é invertível
  \pause
  \[
    \det(A)\det\left(A^{-1}\right)
    \pause
    \;=\; \det\left(AA^{-1}\right)
    \pause
    \;=\; \det(I)
    \pause
    \;=\; 1
  \]

  \pause
  Logo
  \[
    \det(A^{-1})
    \;=\; \frac{1}{\det(A)}
  \]
\end{frame}

%------------------------------------------------------------------------------%
\begin{frame}{Determinante e matriz inversa}
  Para uma matriz quadrada $A$
  \pause
  \[
    \det(A) \neq 0
    \pause
    \qquad\Leftrightarrow\qquad
    A \text{ é invertível}
  \]

  \pause
  Equivalentemente
  \[
    \det(A) = 0
    \pause
    \qquad\Leftrightarrow\qquad
    A \text{ não é invertível}
  \]

  \pause
  Essa é uma das aplicações mais importantes do determinante
\end{frame}

%------------------------------------------------------------------------------%
%------------------------------------------------------------------------------%
\begin{question}
  Verifique se a matriz é invertível
  \[
    A =
    \begin{pmatrix}
      1 & 2 & 3 \\
      0 & 1 & 4 \\
      0 & 0 & 2
    \end{pmatrix}
  \]
\end{question}

%------------------------------------------------------------------------------%
\begin{resolution}{Solução}
  A matriz é triangular
  \pause
  \[
    \det(A)
    \pause
    = a_{11} \; a_{22} \; a_{33}
    \pause
    = 1 \times 1 \times 2
    \pause
    = 2
  \]

  \pause
  Como
  \[
    \det(A) \neq 0
  \]
  \pause
  \(A\) é invertível
\end{resolution}

%------------------------------------------------------------------------------%
\endinput




%------------------------------------------------------------------------------%
\begin{question}
  Verifique se a matriz é invertível
  \[
    B =
    \begin{pmatrix}
      1 & 2 & 3 \\
      2 & 4 & 6 \\
      0 & 1 & 2
    \end{pmatrix}
  \]
\end{question}

%------------------------------------------------------------------------------%
\begin{resolution}{Solução}
  A segunda linha é o dobro da primeira

  \pause
  Portanto
  \[
    \det(B) = 0
  \]

  \pause
  Assim
  \(B\) não é invertível
\end{resolution}

%------------------------------------------------------------------------------%
\endinput

%------------------------------------------------------------------------------%
\begin{question}
  Determine os valores de $k$ para os quais a matriz
  \[
    A =
    \begin{pmatrix}
      1 & 2 & 0 \\
      0 & k & 1 \\
      0 & 0 & 3
    \end{pmatrix}
  \]
  é invertível
\end{question}

%------------------------------------------------------------------------------%
\begin{resolution}{Solução}
  \onslide<+->{Como $A$ é triangular}
  \[
    \onslide<+->{\det(A)}
    \onslide<+->{\;=\; a_{11} \; a_{22} \; a_{33}}
    \onslide<+->{\;=\; 1 \times k \times 3}
    \onslide<+->{\;=\; 3k}
  \]

  \onslide<+->{A matriz é invertível quando}
  \begin{align*}
    \onslide<+->{\det(A) & \neq 0 \\}
    \onslide<+->{3k      & \neq 0 \\}
    \onslide<+->{k       & \neq 0   }
  \end{align*}
\end{resolution}

%------------------------------------------------------------------------------%
\endinput

\endinput
%------------------------------------------------------------------------------%
\begin{question}
\end{question}

%------------------------------------------------------------------------------%
\begin{resolution}{Solução}
\end{resolution}

%------------------------------------------------------------------------------%
\endinput

\begin{frame}{Exemplo -- Laplace em uma matriz $4\times4$}
Calcule
\[
\det(A),
\qquad
A=
\begin{pmatrix}
1&0&0&2\\
0&3&0&0\\
4&0&5&0\\
0&0&0&2
\end{pmatrix}.
\]
\medskip

A primeira coluna possui dois zeros, mas a quarta coluna possui
três zeros. Expandimos pela quarta coluna:
\[
\det(A)
=
2C_{44}.
\]
\medskip

Como $4+4$ é par,
\[
C_{44}=M_{44}.
\]
\medskip

Logo,
\[
\det(A)
=
2
\begin{vmatrix}
1&0&0\\
0&3&0\\
4&0&5
\end{vmatrix}.
\]
\end{frame}

\begin{frame}{Exemplo -- Laplace em uma matriz $4\times4$}
A matriz restante é triangular inferior:
\[
\begin{pmatrix}
1&0&0\\
0&3&0\\
4&0&5
\end{pmatrix}.
\]
Portanto,
\[
M_{44}=1\cdot3\cdot5=15.
\]
Assim,
\[
\det(A)=2\cdot15=30.
\]
\medskip

Logo,
\[
\boxed{\det(A)=30}.
\]
\end{frame}

\endinput
%------------------------------------------------------------------------------%
\begin{question}
\end{question}

%------------------------------------------------------------------------------%
\begin{resolution}{Solução}
\end{resolution}

%------------------------------------------------------------------------------%
\endinput

\begin{frame}{Exemplo -- Escolhendo a melhor expansão}
Considere
\[
A=
\begin{pmatrix}
1&2&0&0\\
0&3&0&0\\
2&1&4&5\\
0&0&6&7
\end{pmatrix}.
\]
\medskip

A segunda coluna não possui muitos zeros, enquanto a segunda
linha possui três zeros. Portanto, é conveniente expandir pela
segunda linha:
\[
\det(A)
=
3C_{22}.
\]
\medskip

Como $2+2$ é par,
\[
C_{22}=M_{22}.
\]
\end{frame}

\begin{frame}{Exemplo -- Escolhendo a melhor expansão}
Temos
\[
M_{22}
=
\begin{vmatrix}
1&0&0\\
2&4&5\\
0&6&7
\end{vmatrix}.
\]
\medskip

Expandindo pela primeira linha:
\[
M_{22}
=
\begin{vmatrix}
4&5\\
6&7
\end{vmatrix}
=28-30=-2.
\]
\medskip

Portanto,
\[
\det(A)=3(-2)=-6.
\]
\medskip

Assim,
\[
\boxed{\det(A)=-6}.
\]
\end{frame}



%------------------------------------------------------------------------------%
\section{Resumo}

%------------------------------------------------------------------------------%
\begin{frame}{Resumo das principais propriedades}
  Para uma matriz quadrada $A$
  \begin{enumerate}
    \item matriz triangular
    \(
      \det(A) = a_{11} a_{22} \cdots a_{nn}
    \)
    \item troca de duas linhas muda o sinal do determinante
    \item multiplicação de uma linha por $k$ multiplica o determinante por $k$
    \item adicionar um múltiplo de uma linha a outra não altera o determinante
    \item uma linha nula implica determinante nulo
    \item duas linhas iguais ou proporcionais implicam determinante nulo
    \item \(\det(A^t) = \det(A)\)
    \item \(\det(AB) = \det(A)\det(B)\)
  \end{enumerate}
\end{frame}

%------------------------------------------------------------------------------%
% \framedgraphic{Gráfico de uma Função Real}{B_01_grafico_funcao_real}{}

%------------------------------------------------------------------------------%
%\section{Lista Mínima}
%
%%------------------------------------------------------------------------------%
%\begin{frame}{Lista Mínima}
%  \begin{enumerate}
%    \item
%  \end{enumerate}
%
%  \vfill
%  Atenção: A prova é baseada no livro, não nas apresentações
%\end{frame}

%------------------------------------------------------------------------------%
\end{document}
%------------------------------------------------------------------------------%
